\subsection{Four under-specifications of the reference}
\label{sec:nulls}

A content floor is not a dataset constant. Four choices usually left implicit each move it by more than the effects under study, so all four are part of the measured construct and must be printed with the score. \emph{(i) The tie convention:} on MMLU-Pro, one dataset and one tokenizer give 0.125914 under \texttt{wrong} and 0.532164 under \texttt{credit}, a 40.6-point span whose upper end exceeds the intact base model's own \texttt{content\_norm} score by 32.5 points, and on ARC-Challenge \texttt{credit} alone moves five of six OLMo-2 arms from significantly above the floor to significantly below it. \emph{(ii) The length unit}, characters or continuation tokens: OpenBookQA \texttt{credit} is 0.416 versus 0.644, ARC-Challenge 18.60\% tied-longest under characters against 50.85\% under tokens. \emph{(iii) The tokenizer}, which moves \texttt{credit} by up to 10.6 points on four-way tasks and 9.26 on MMLU-Pro, non-monotonically in vocabulary size \citep{oostermeijer2026length}. \emph{(iv) The meaning of ``chance''} when \texttt{n\_opt} runs from three to ten: naive $1/10$ gives 0.100000 and $\mathrm{mean}(1/\texttt{n\_opt})$ gives 0.110877, so the gap to the 0.116606 always-A floor is either $+1.661$ points and $1.1661\times$ or $+0.573$ points and $1.0517\times$; both must be reported. The four are independent, so a token content floor is fixed only once all four are --- a property of $(\text{dataset},\text{convention},\text{unit},\text{tokenizer})$, not of the dataset. The letter floor is exempt from all four: a pure property of the item set's gold labels, invariant across all 15 cross-family arms and bit-identical across all 21 MMLU-Pro cells. Appendix~\ref{app:nulls} develops each with Tables~\ref{tab:nulls} and~\ref{tab:conventions}.
